\documentclass[12pt]{article}
\usepackage[T1]{fontenc}
\usepackage[utf8]{inputenc}
\usepackage{lmodern}
\usepackage{textcomp}
\usepackage{enumitem}
\usepackage{geometry}
\usepackage{graphicx}
\usepackage{amsmath, amssymb}
\geometry{margin=1in}
\begin{document}
\begin{center}
Constitutional Law - Undergraduate Level Assessment 4 \
Total Marks: 40 \
Answer all questions.
\end{center}

\section*{Multiple Choice (10 marks)}
\begin{enumerate}[label=\arabic*.]
\item Which statement below is the least likely to follow logically from Dworkin's notion of law as integrity?
\begin{enumerate}[label=(\alph*)]
    \item It is likely to generate more individual rights and greater liberty.
    \item It renders a community more genuine.
    \item It opens the door to authoritarianism.
    \item It improves the moral justification for the exercise of political power.
\end{enumerate}
\item The U.S. Congress is empowered by the Commerce Clause and other provisions of the U.S. Constitution to enact \_\_\_\_\_\_\_\_ to regulate foreign and interstate commerce.
\begin{enumerate}[label=(\alph*)]
    \item ordinances
    \item federal statutes
    \item executive orders
    \item charters
\end{enumerate}
\item What is the purpose of sovereign immunity?
\begin{enumerate}[label=(\alph*)]
    \item The purpose of immunity is to protect foreign Heads of State from embarrassment
    \item Immunity protects a State from being invaded by another
    \item Immunity shields States from being sued in the courts of other States
    \item The purpose of immunity is to offer impunity in respect of all crimes
\end{enumerate}
\item Why does Dworkin support liberal egalitarianism?
\begin{enumerate}[label=(\alph*)]
    \item Because it attempts to give effect to personal choice over individual luck.
    \item Because liberty is more important than equality.
    \item Because a market economy is just.
    \item Because the state is the best arbiter of equality between individuals.
\end{enumerate}
\item Is recognition of governments prevalent in contemporary international practice?
\begin{enumerate}[label=(\alph*)]
    \item Recognition of governments is very prevalent in contemporary practice
    \item Recognition of governments has largely been replaced by functional Recognition
    \item Government recognition is common in respect of rebel entities
    \item Only democratic governments are recognised in contemporary practice
\end{enumerate}
\end{enumerate}

\section*{True / False (10 marks)}
\textit{Answer True or False}
\begin{enumerate}[label=\arabic*.]
\setcounter{enumi}{5}
\item True or False: Nozick's formula r x H represents the effectiveness of rehabilitation multiplied by the hazard to the community.
\item True or False: Hume contends that natural law is primarily concerned with historical precedents rather than moral objectivity.
\item True or False: An 'armed attack' in the UN Charter allows for the invasion of any State regardless of the intensity of the armed force used.
\item True or False: The Badinter Commission required a commitment to human rights and democracy from the former Yugoslav republics.
\item True or False: The UN Human Rights Committee is a treaty-based mechanism that monitors compliance with the International Covenant on Civil and Political Rights.
\end{enumerate}

\section*{Long Form Response (20 marks)}
\textit{Answer each question with a clear and complete explanation.}
\begin{enumerate}[label=\arabic*.]
\setcounter{enumi}{10}
\item Discuss Durkheim's perspective on the function of punishment within society, particularly focusing on the role of vengeance and its implications for social cohesion and justice.
\item Discuss John Finnis's critique of Hume's conception of practical reason, focusing on how Finnis argues that human reason aids in determining a worthwhile life.
\end{enumerate}

\vfill
\noindent\textit{End of Paper}
\end{document}